%-------------------------
% Shared preamble for resume.tex and cv.tex
% Template: Jake's Resume, MIT licensed, https://github.com/jakegut/resume
% Based off of: https://github.com/sb2nov/resume
%------------------------

% \iffullcv switches content.tex between the full CV (true) and the resume (false)
\newif\iffullcv

\usepackage{latexsym}
\usepackage[empty]{fullpage}
\usepackage{titlesec}
\usepackage{marvosym}
\usepackage[usenames,dvipsnames]{color}
\usepackage{verbatim}
\usepackage{enumitem}
\usepackage[hidelinks]{hyperref}
\usepackage{fancyhdr}
\usepackage[english]{babel}
\usepackage{tabularx}
\input{glyphtounicode}

%-------------------------
% Bibliography: ACS-style entries from the website's bib files (one source for
% the site, the CV and the resume). Paths are relative to this cv/ folder.
\usepackage[backend=biber, style=chem-acs, articletitle=true, doi=true, url=false,
            maxnames=99, sorting=none, eventdate=comp, defernumbers=true]{biblatex}
\addbibresource{../_bibliography/publications.bib}
\addbibresource{../_bibliography/presentations.bib}
\DeclareSourcemap{\maps[datatype=bibtex]{%
  % Print the abbreviated journal name (shortjournal) when the entry has one
  \map[overwrite]{%
    \step[fieldsource=shortjournal]\step[fieldset=journaltitle, origfieldval]}%
  % Tag every entry from presentations.bib with the keyword "presentation"
  % (appended, so any keywords already in the entry are kept)
  \map[overwrite]{\perdatasource{../_bibliography/presentations.bib}%
    \step[fieldset=keywords, fieldvalue={,}, appendstrict]%
    \step[fieldset=keywords, fieldvalue=presentation, append]}%
  % Entries with selected = {true} (the same flag the website uses for its home
  % page) get the keyword "selected"; the resume prints only those
  \map[overwrite]{%
    \step[fieldsource=selected, match=\regexp{\A\s*true\s*\Z}, final]%
    \step[fieldset=keywords, fieldvalue={,}, appendstrict]%
    \step[fieldset=keywords, fieldvalue=selected, append]}%
}}
\defbibfilter{publications}{not keyword=presentation and not keyword=preprint}
\defbibfilter{preprints}{keyword=preprint}
\defbibfilter{presentations}{keyword=presentation}
\defbibfilter{selectedpublications}{not keyword=presentation and keyword=selected}
\defbibfilter{selectedpresentations}{keyword=presentation and keyword=selected}
% Meeting presentations (@unpublished), ACS format:
%   Authors. Title. Presented at Meeting, Venue, Date; type number. Award.
\DeclareBibliographyDriver{unpublished}{%
  \usebibmacro{bibindex}%
  \usebibmacro{begentry}%
  \usebibmacro{author/translator+others}%
  \setunit{\labelnamepunct}\newblock
  \usebibmacro{title}%
  \setunit{\adddot\space}\newblock
  \iffieldundef{eventtitle}{}{\printtext{Presented at}\setunit{\addspace}}%
  \printfield{eventtitle}%
  \setunit{\addcomma\space}%
  \printfield{venue}%
  \setunit{\addcomma\space}%
  \printeventdate
  \setunit{\addsemicolon\space}%
  \printfield{type}%
  \setunit{\addspace}%
  \printfield{number}%
  \setunit{\adddot\space}\newblock
  \usebibmacro{addendum+pubstate}%
  \usebibmacro{finentry}}
% Bold the author marked with author+an = {N=highlight} in the .bib file,
% and add * after authors marked "corresponding" (e.g. {2=highlight,corresponding})
\renewcommand*{\mkbibnamegiven}[1]{\ifitemannotation{highlight}{\textbf{#1}}{#1}}
\renewcommand*{\mkbibnamefamily}[1]{\ifitemannotation{highlight}{\textbf{#1}}{#1}}
\renewcommand*{\mkbibcompletename}[1]{#1\ifitemannotation{corresponding}{\textsuperscript{*}}{}}
% Unnumbered list matching the other CV sections
\defbibenvironment{cvpubs}
  {\begin{itemize}[leftmargin=0.15in, label={}, itemsep=2pt]\small}
  {\end{itemize}}
  {\item}

\pagestyle{fancy}
\fancyhf{} % clear all header and footer fields
\fancyfoot{}
\renewcommand{\headrulewidth}{0pt}
\renewcommand{\footrulewidth}{0pt}

% Adjust margins
\addtolength{\oddsidemargin}{-0.5in}
\addtolength{\evensidemargin}{-0.5in}
\addtolength{\textwidth}{1in}
\addtolength{\topmargin}{-.5in}
\addtolength{\textheight}{1.0in}

\urlstyle{same}

\raggedbottom
\raggedright
\setlength{\tabcolsep}{0in}

% Sections formatting
\titleformat{\section}{
  \vspace{-4pt}\scshape\raggedright\large
}{}{0em}{}[\color{black}\titlerule \vspace{-5pt}]

% Ensure that generate pdf is machine readable/ATS parsable
\pdfgentounicode=1

%-------------------------
% Custom commands
\newcommand{\resumeItem}[1]{
  \item\small{
    {#1 \vspace{-2pt}}
  }
}

\newcommand{\resumeSubheading}[4]{
  \vspace{-2pt}\item
    \begin{tabular*}{0.97\textwidth}[t]{l@{\extracolsep{\fill}}r}
      \textbf{#1} & #2 \\
      \textit{\small#3} & \textit{\small #4} \\
    \end{tabular*}\vspace{-7pt}
}

\newcommand{\resumeProjectHeading}[2]{
    \item
    \begin{tabular*}{0.97\textwidth}{l@{\extracolsep{\fill}}r}
      \small#1 & #2 \\
    \end{tabular*}\vspace{-7pt}
}

% Right-aligned trailer (date) that never splits: if it does not fit on the
% last line it moves, whole and still right-aligned, to a new line (TeXbook trick)
\newcommand{\cvright}[1]{\unskip\nobreak\hfill\penalty50\hskip1em\hbox{}\nobreak\hfill\mbox{#1}}

\newcommand{\resumeSubItem}[1]{\resumeItem{#1}\vspace{-4pt}}

\renewcommand\labelitemii{$\vcenter{\hbox{\tiny$\bullet$}}$}

\newcommand{\resumeSubHeadingListStart}{\begin{itemize}[leftmargin=0.15in, label={}]}
\newcommand{\resumeSubHeadingListEnd}{\end{itemize}}
\newcommand{\resumeItemListStart}{\begin{itemize}}
\newcommand{\resumeItemListEnd}{\end{itemize}\vspace{-5pt}}
